\uuid{kFvZ}
\titre{Espaces supplémentaires/Base/Dimension}
\auteur{hanani ; mzari}
\organisation{exo7}
\licence{CC BY-NC-SA 4.0}
\datecreate{2020-01-20}
\niveau{L1}
\module{Algèbre}
\chapitre{Espace vectoriel}
\sousChapitre{Somme directe}
\difficulte{1}
\theme{Espaces vectoriels supplémentaires, Espaces supplémentaires, Base, Dimension}
\exo7qcmid{lille-2.espaces-vectoriels.051}

\contenu{
\begin{question}
\qclassification{106.03}
\qtags{motcle=Espaces supplémentaires/Base/Dimension}
On considère les deux sous-espaces vectoriels de $\Rr^4$ :
$$E= \mbox {vect} \{u_1,u_2,u_3\},\mbox{  où }u_1=(1,-1,0,1), \; u_2=(1,0,1,0),\; u_3=(3,-1,1,2)$$
et
$$F=\{(x,y,z,t)\in \Rr^4\, ;\, x+y-z=0\; \mbox{et}\; y+z=0\}.$$
Quelles sont les assertions vraies ?
\begin{answers}
\good{$\dim E = 3$}.
\good{$\dim E\cap F = 1$}.
\good{$E+F= \Rr^4$}.
\bad{$E$ et $F$ sont supplémentaires dans $\Rr^4$}.
\end{answers}
\begin{explanations} On vérifie que $\{u_1,u_2,u_3\}$ est libre  et
que $F=\mbox {vect} \{v_1,v_2\}$, où $v_1=(2,-1,1,0)$ et $ v_2=(0,0,0,1)$. Par conséquent, $\dim E=3$ et $\dim F=2$.
\vskip1mm
Il y a une seule relation de dépendance entre $u_1,u_2,u_3,v_1$ et $v_2$. Soit : $u_1+u_2-v_1-v_2=0$. On déduit que $\{u_1,u_2,u_3, v_1\}$ est une base de $E+F$, donc $\dim (E+F)=4$ et comme $E+F$ est un sous-espace de $\Rr^4$ et $\dim \Rr^4=4$, $E+F= \Rr^4$. Du théorème de la dimension d'une somme, on déduit que $\dim E\cap F=1$, donc $E$ et $F$ ne sont pas supplémentaires dans $\Rr^4$.
\end{explanations}
\end{question}
}
